\chapter{Turunan: Perkenalan Pertama}\label{ch:g11:deriv}

Seberapa cepat sebuah \omterm{def:g11:func:function}{fungsi} berubah? Jawabannya di satu titik adalah
\emph{\omterm{def:g11:deriv:derivative}{turunan}}, bilangan tunggal yang paling berguna dalam seluruh
matematika terapan. Bab ini membangunnya dari \omterm{def:g11:deriv:rate}{laju perubahan rata-rata},
menghitungnya untuk \omterm{def:g11:func:reference}{fungsi acuan}, lalu memakai tandanya untuk membaca
variasi sebuah \omterm{def:g11:func:function}{fungsi}. Semua yang ada di sini diperluas dan diperdalam pada
\cref{ch:g12:deriv}.

\section{Laju perubahan rata-rata}

\begin{definition}[Laju perubahan rata-rata]\label{def:g11:deriv:rate}
Misalkan $f$ terdefinisi pada \omterm{def:g10:numbers:interval}{interval} $I$, serta $a \neq b$ di $I$.
\emph{Laju perubahan rata-rata}\index{laju perubahan rata-rata} $f$ antara
$a$ dan $b$ adalah
\[
\frac{f(b) - f(a)}{b - a}.
\]
Bilangan itu adalah \omterm{def:g10:reffunc:affine}{gradien} \emph{garis sekan}\index{garis sekan} yang melalui
titik $\bigl(a, f(a)\bigr)$ dan $\bigl(b, f(b)\bigr)$ pada kurvanya.
\end{definition}

\begin{example}
Sebuah mobil menempuh $140$ km antara pukul 14.00 dan 16.00: laju rata-ratanya
$\frac{140}{2} = 70$ km/jam. Namun spidometernya menampilkan
laju \emph{sesaat} pada setiap saat --- dan justru itulah gagasan
yang dipertepat oleh \omterm{def:g11:deriv:derivative}{turunan}.
\end{example}

\section{Turunan di sebuah titik}

\begin{definition}[Turunan di sebuah titik]\label{def:g11:deriv:derivative}
Misalkan $f$ terdefinisi pada \omterm{def:g10:numbers:interval}{interval} yang memuat $a$. Untuk $h \neq 0$,
bentuklah \omterm{def:g11:deriv:rate}{laju perubahan rata-rata} antara $a$ dan $a + h$:
\[
r(h) = \frac{f(a+h) - f(a)}{h}.
\]
Jika $r(h)$ mendekati satu bilangan tetap ketika $h$ menjadi sedekat mungkin ke
$0$, kita katakan bahwa $f$ \emph{dapat diturunkan di $a$}\index{turunan}\index{dapat diturunkan};
bilangan itu disebut \emph{turunan $f$ di $a$}, ditulis $f'(a)$:
\[
f'(a) = \lim_{h \to 0} \frac{f(a+h) - f(a)}{h}.
\]
\end{definition}

\begin{example}\label{ex:g11:deriv:atpoint}
Misalkan $f(x) = x^2$ dan $a = 1$. Untuk $h \neq 0$:
\[
\frac{f(1+h) - f(1)}{h}
= \frac{(1+h)^2 - 1}{h}
= \frac{1 + 2h + h^2 - 1}{h}
= \frac{h(2 + h)}{h}
= 2 + h .
\]
Ketika $h$ mendekati $0$, nilai ini mendekati $2$: \omterm{def:g11:func:function}{fungsi} $x^2$
\omterm{def:g11:deriv:derivative}{dapat diturunkan} di $1$ dan $f'(1) = 2$. Perhatikan siasatnya:
\emph{sederhanakan hasil baginya sampai $h$ tidak lagi muncul pada penyebut},
lalu biarkan $h \to 0$.
\end{example}

\begin{example}[Sebuah titik yang tak dapat diturunkan]
Misalkan $f(x) = \abs{x}$ dan $a = 0$. Maka
$\frac{f(h) - f(0)}{h} = \frac{\abs h}{h}$, yang bernilai $1$ untuk $h > 0$
dan $-1$ untuk $h < 0$: tidak ada satu bilangan pun yang didekati, jadi
$\abs x$ tidak \omterm{def:g11:deriv:derivative}{dapat diturunkan} di $0$. Kurvanya mempunyai \emph{patahan} di sana.
\end{example}

\section{Garis singgung}

\begin{definition}[Garis singgung]\label{def:g11:deriv:tangent}
Jika $f$ \omterm{def:g11:deriv:derivative}{dapat diturunkan} di $a$, \emph{garis singgung}\index{garis singgung}
kurvanya di titik $A = \bigl(a, f(a)\bigr)$ adalah garis yang melalui $A$
dengan \omterm{def:g10:reffunc:affine}{gradien} $f'(a)$. \omterm{def:g10:algebra:equation}{Persamaannya} adalah
\[
y = f(a) + f'(a)(x - a).
\]
\end{definition}

\omterm{def:g11:deriv:tangent}{Garis singgung} adalah kedudukan batas \omterm{def:g11:deriv:rate}{garis sekan} yang melalui $A$: ketika
$h \to 0$, titik keduanya $\bigl(a+h, f(a+h)\bigr)$ meluncur sepanjang
kurvanya menuju $A$, dan \omterm{def:g10:reffunc:affine}{gradien} sekannya $\frac{f(a+h)-f(a)}{h}$
mendekati $f'(a)$.

\begin{omfigure}
\begin{tikzpicture}
\begin{axis}[omaxis,
  width=8.8cm, height=6cm,
  xmin=-0.4, xmax=3.2, ymin=-0.6, ymax=4.3,
  xtick={1,2.2}, xticklabels={$a$,$a+h$}, ytick=\empty]
  \addplot[omProp, thick, domain=-0.1:2.9] {1 + 1.6*(x-1)};
  \addplot[omThm, thick, domain=-0.35:2.6] {1 + 1*(x-1)};
  \addplot[omDef, thick, domain=-0.4:2.75] {x^2/2 + 0.5};
  \fill (axis cs:1,1) circle (1.5pt) node[above left, font=\small] {$A$};
  \fill (axis cs:2.2,2.92) circle (1.5pt)
    node[above left, font=\small] {$B$};
  \node[omProp, font=\small, anchor=west] at (axis cs:2.62,3.75)
    {sekan};
  \node[omThm, font=\small, anchor=west] at (axis cs:2.62,2.45)
    {singgung};
\end{axis}
\end{tikzpicture}

{\small \omterm{def:g11:deriv:rate}{Garis sekan} yang melalui $A$ dan $B = \bigl(a+h, f(a+h)\bigr)$
bergradien $\frac{f(a+h)-f(a)}{h}$. Ketika $B$ meluncur menuju $A$, sekannya
memiring menjadi \omterm{def:g11:deriv:tangent}{garis singgung}, yang bergradien $f'(a)$.}
\end{omfigure}

\begin{example}\label{ex:g11:deriv:tangenteq}
Untuk $f(x) = x^2$ di $a = 1$: $f(1) = 1$ dan $f'(1) = 2$
(\cref{ex:g11:deriv:atpoint}), sehingga \omterm{def:g11:deriv:tangent}{garis singgungnya} adalah
\[
y = 1 + 2(x - 1) = 2x - 1 .
\]
Inilah garis $d_2$ yang ditemukan lewat aljabar murni pada \cref{exo:g11:quad:11}.
\end{example}

\section{Turunan fungsi acuan}

\begin{proposition}\label{prop:g11:deriv:refderivatives}
Pada \omterm{def:g11:func:function}{daerah asal} tempat \omterm{def:g11:func:function}{fungsi} itu \omterm{def:g11:deriv:derivative}{dapat diturunkan}:
\[
(c)' = 0, \quad
(x)' = 1, \quad
(x^2)' = 2x, \quad
(x^3)' = 3x^2, \quad
\left(\frac1x\right)' = -\frac{1}{x^2}, \quad
(\sqrt x\,)' = \frac{1}{2\sqrt x} \ (x > 0).
\]
Secara lebih umum, $(x^n)' = n\,x^{n-1}$ untuk setiap \omterm{def:g10:numbers:sets}{bilangan bulat} $n \geq 1$.
\end{proposition}

\begin{proof}
Setiap buktinya mengikuti siasat \cref{ex:g11:deriv:atpoint}: sederhanakan
lajunya, lalu biarkan $h \to 0$. Tetapkan $a$ di \omterm{def:g11:func:function}{daerah asalnya}.

\emph{\omterm{def:g11:func:function}{Fungsi} tetap dan \omterm{def:g11:func:function}{fungsi} identitas.} $\frac{c - c}{h} = 0$ dan
$\frac{(a+h) - a}{h} = 1$ untuk semua $h$: limitnya $0$ dan $1$.

\emph{Pangkat dua.} $\frac{(a+h)^2 - a^2}{h} = \frac{2ah + h^2}{h} = 2a + h
\to 2a$.

\emph{Pangkat tiga.} $(a+h)^3 = a^3 + 3a^2h + 3ah^2 + h^3$, sehingga lajunya
$3a^2 + 3ah + h^2 \to 3a^2$.

\emph{Kebalikan.} Untuk $a \neq 0$ dan $h$ yang cukup kecil sehingga
$a + h \neq 0$:
\[
\frac{1}{h}\left(\frac{1}{a+h} - \frac{1}{a}\right)
= \frac{1}{h} \cdot \frac{a - (a+h)}{a(a+h)}
= \frac{-1}{a(a+h)}
\xrightarrow[h \to 0]{} -\frac{1}{a^2}.
\]

\emph{\omterm{def:g11:quad:discriminant}{Akar} kuadrat.} Untuk $a > 0$, kalikan dengan bentuk sekawannya:
\[
\frac{\sqrt{a+h} - \sqrt a}{h}
= \frac{(a + h) - a}{h\,(\sqrt{a+h} + \sqrt a)}
= \frac{1}{\sqrt{a+h} + \sqrt a}
\xrightarrow[h \to 0]{} \frac{1}{2\sqrt a}.
\]

\emph{Pangkat umum.} Polanya, $2a$, $3a^2$, berlanjut: setelah
$(a+h)^n$ dijabarkan, suku-suku setelah $a^n + n\,a^{n-1}h$ semuanya membawa
$h^2$, sehingga lajunya menuju $n\,a^{n-1}$. (Induksi lengkapnya diberikan pada
\cref{ch:g12:deriv}.)
\end{proof}

\section{Aturan penurunan}

\begin{theorem}[Operasi]\label{thm:g11:deriv:rules}
Misalkan $u$ dan $v$ \omterm{def:g11:deriv:derivative}{dapat diturunkan} pada $I$ dan $\lambda \in \R$. Maka
$u + v$, $\lambda u$, $uv$ \omterm{def:g11:deriv:derivative}{dapat diturunkan} pada $I$, demikian pula
$\frac{u}{v}$ di mana pun $v \neq 0$, dan
\[
(u+v)' = u' + v', \qquad
(\lambda u)' = \lambda u', \qquad
(uv)' = u'v + uv', \qquad
\left(\frac{u}{v}\right)' = \frac{u'v - uv'}{v^2}.
\]
\end{theorem}

\begin{proof}[Bukti tiga aturan pertama]
Tetapkan $a \in I$ lalu singkat $\Delta u = u(a+h) - u(a)$,
$\Delta v = v(a+h) - v(a)$.

\emph{Jumlah.} Laju $u + v$ adalah
$\frac{\Delta u + \Delta v}{h} = \frac{\Delta u}{h} + \frac{\Delta v}{h}
\to u'(a) + v'(a)$.

\emph{Kelipatan.} Laju $\lambda u$ adalah
$\lambda \frac{\Delta u}{h} \to \lambda u'(a)$.

\emph{Hasil kali.} Tambahkan lalu kurangkan suku campuran $u(a+h)\,v(a)$:
\begin{align*}
\frac{u(a+h)v(a+h) - u(a)v(a)}{h}
&= \frac{u(a+h)v(a+h) - u(a+h)v(a) + u(a+h)v(a) - u(a)v(a)}{h}\\
&= u(a+h)\,\frac{\Delta v}{h} + v(a)\,\frac{\Delta u}{h}.
\end{align*}
Ketika $h \to 0$: $\frac{\Delta v}{h} \to v'(a)$, $\frac{\Delta u}{h} \to
u'(a)$, dan $u(a+h) \to u(a)$ (lajunya mempunyai limit, sehingga
$\Delta u = h \cdot \frac{\Delta u}{h} \to 0$). Jadi lajunya menuju
$u(a)v'(a) + v(a)u'(a)$.

Aturan hasil bagi dibuktikan menurut garis yang sama; aturan itu
\emph{diterima tanpa bukti pada tingkat ini} dan dibuktikan pada \cref{ch:g12:deriv}.
\end{proof}

\begin{example}\label{ex:g11:deriv:computations}
\emph{Polinomial.} $f(x) = 3x^3 - 5x^2 + 7x - 2$:
\[
f'(x) = 3 \times 3x^2 - 5 \times 2x + 7 = 9x^2 - 10x + 7 .
\]

\emph{Hasil bagi.} $g(x) = \frac{2x+1}{x^2+1}$: dengan $u = 2x + 1$, $v = x^2
+ 1$,
\[
g'(x) = \frac{2(x^2+1) - (2x+1)\,2x}{(x^2+1)^2}
= \frac{-2x^2 - 2x + 2}{(x^2+1)^2}.
\]

\emph{Hasil kali.} $k(x) = x\sqrt x$ untuk $x > 0$:
$k'(x) = 1 \times \sqrt x + x \times \frac{1}{2\sqrt x}
= \sqrt x + \frac{\sqrt x}{2} = \frac{3\sqrt x}{2}$.
\end{example}

\section{Tanda turunan dan variasi}

\begin{theorem}[Turunan dan variasi]\label{thm:g11:deriv:variations}
Misalkan $f$ \omterm{def:g11:deriv:derivative}{dapat diturunkan} pada sebuah \omterm{def:g10:numbers:interval}{interval} $I$.
\begin{enumerate}
  \item Jika $f' \geq 0$ pada $I$, maka $f$ naik pada $I$.
  \item Jika $f' \leq 0$ pada $I$, maka $f$ turun pada $I$.
  \item Jika $f' > 0$ pada $I$ (kecuali mungkin di berhingga banyak titik
        tempat \omterm{def:g11:deriv:derivative}{turunannya} lenyap), maka $f$ tegas naik pada $I$.
\end{enumerate}
\end{theorem}

\begin{proof}
Gagasannya jelas --- kurva yang semua \omterm{def:g11:deriv:tangent}{garis singgungnya} menunjuk ke atas pasti
mendaki --- tetapi bukti yang cermat memerlukan teorema nilai rata-rata. Hasil
ini \emph{diterima tanpa bukti pada tingkat ini}; lihat \cref{thm:g12:deriv:variations}.
\end{proof}

\begin{proposition}[Nilai ekstrem]\label{prop:g11:deriv:extrema}
Jika $f'$ berganti tanda di $a$ dari $+$ ke $-$, maka $f$ mempunyai \omterm{def:g10:functions:extrema}{maksimum}
lokal di $a$; dari $-$ ke $+$, \omterm{def:g10:functions:extrema}{minimum} lokal. Di titik semacam itu $f'(a) = 0$
dan \omterm{def:g11:deriv:tangent}{garis singgungnya} mendatar.
\end{proposition}

\begin{proof}
Jika $f' \geq 0$ pada \omterm{def:g10:numbers:interval}{interval} yang berakhir di $a$ dan $f' \leq 0$ pada
\omterm{def:g10:numbers:interval}{interval} yang bermula di $a$, maka $f$ naik sampai $f(a)$ lalu turun sesudahnya
(\cref{thm:g11:deriv:variations}): jadi $f(a)$ terbesar di sekitarnya. Karena
$f'$ berganti tanda di $a$ dan terdefinisi di sana, $f'(a) = 0$.
\end{proof}

\begin{method}[Menelaah sebuah fungsi]\label{met:g11:deriv:study}
\begin{enumerate}
  \item Hitunglah $f'(x)$ lalu sederhanakan --- sebaiknya menjadi bentuk terfaktorkan.
  \item Tentukan tanda $f'(x)$ pada setiap \omterm{def:g10:numbers:interval}{interval} (untuk $f'$ yang
        berbentuk kuadrat, pakailah \cref{thm:g11:quad:sign}).
  \item Catat hasilnya dalam sebuah \emph{\omterm{def:g10:functions:table}{tabel variasi}}: satu baris untuk
        tanda $f'$, satu baris panah untuk $f$, beserta nilai $f$ di
        titik-titik baliknya.
  \item Simpulkan: nilai ekstrem, banyaknya penyelesaian $f(x) = k$, ketaksamaan.
\end{enumerate}
\end{method}

\begin{example}\label{ex:g11:deriv:study}
Telaahlah $f(x) = x^3 - 3x + 1$ pada $\R$. Pertama,
$f'(x) = 3x^2 - 3 = 3(x-1)(x+1)$: sebuah bentuk kuadrat dengan \omterm{def:g11:quad:discriminant}{akar} $-1$ dan
$1$, yang positif di luar keduanya. Jadi $f$ naik pada $\intoc{-\infty}{-1}$,
turun pada $\intcc{-1}{1}$, naik pada $\intco{1}{+\infty}$, dengan
\omterm{def:g10:functions:extrema}{maksimum} lokal $f(-1) = -1 + 3 + 1 = 3$ dan \omterm{def:g10:functions:extrema}{minimum} lokal
$f(1) = 1 - 3 + 1 = -1$.
\end{example}

\begin{omfigure}
\begin{tikzpicture}
\begin{axis}[omaxis,
  width=8.4cm, height=5.8cm,
  xmin=-2.6, xmax=2.6, ymin=-3.6, ymax=4.4,
  xtick={-1,1}, ytick={-1,3}]
  \addplot[omProp, thick, domain=-1.7:-0.3] {3};
  \addplot[omProp, thick, domain=0.3:1.7] {-1};
  \addplot[omDef, thick, domain=-2.25:2.25, samples=90] {x^3 - 3*x + 1};
  \fill (axis cs:-1,3) circle (1.5pt);
  \fill (axis cs:1,-1) circle (1.5pt);
\end{axis}
\end{tikzpicture}

{\small Kurva $f(x) = x^3 - 3x + 1$
(\cref{ex:g11:deriv:study}): \omterm{def:g11:deriv:tangent}{garis singgung} mendatar (jingga) di
\omterm{def:g10:functions:extrema}{maksimum} lokal $(-1, 3)$ dan \omterm{def:g10:functions:extrema}{minimum} lokal $(1, -1)$, tempat $f'$ berganti
tanda.}
\end{omfigure}

\begin{example}[Pengoptimuman]\label{ex:g11:deriv:box}
Sebuah kotak terbuka dibuat dari lembaran $12 \times 12$ cm dengan memotong
persegi bersisi $x$ dari setiap pojoknya lalu melipatnya. Volumenya adalah
$V(x) = x(12 - 2x)^2$ untuk $0 < x < 6$. Setelah dijabarkan,
$V(x) = 4x^3 - 48x^2 + 144x$, sehingga
\[
V'(x) = 12x^2 - 96x + 144 = 12(x^2 - 8x + 12) = 12(x - 2)(x - 6).
\]
Pada $\intoo{0}{6}$, faktor $x - 6$ bernilai negatif, sehingga $V'(x)$ bertanda
sama dengan $-(x-2)$: positif sebelum $2$, negatif sesudahnya. Volumenya \omterm{def:g10:functions:extrema}{maksimum}
untuk $x = 2$ cm, yaitu $V(2) = 2 \times 8^2 = 128$ cm$^3$.
\end{example}

\section{Latihan}

\begin{exercise}[$\star$]\label{exo:g11:deriv:1}
Dengan memakai definisinya (laju perubahan, lalu $h \to 0$), hitunglah $f'(2)$
untuk $f(x) = x^2$, lalu $g'(1)$ untuk $g(x) = x^2 + 3x$.
\end{exercise}

\begin{exercise}[$\star$]\label{exo:g11:deriv:2}
Turunkan:
\[
f(x) = 4x^3 - 2x^2 + x - 7, \qquad
g(x) = 5\sqrt{x} + \frac{3}{x}, \qquad
h(x) = (2x+1)(x^2 - 3) .
\]
\end{exercise}

\begin{exercise}[$\star$]\label{exo:g11:deriv:3}
Berikan \omterm{def:g10:algebra:equation}{persamaan} \omterm{def:g11:deriv:tangent}{garis singgung} kurva $f$ di titik yang
berabsis $a$:
\[
f(x) = x^2 - 3x, \ a = 2; \qquad
f(x) = \frac1x, \ a = 1; \qquad
f(x) = \sqrt x, \ a = 4 .
\]
\end{exercise}

\begin{exercise}[$\star$]\label{exo:g11:deriv:4}
Turunkan hasil bagi berikut:
\[
f(x) = \frac{x+2}{x-1}, \qquad
g(x) = \frac{x^2}{x^2+1}, \qquad
h(x) = \frac{1}{x^2 + x + 1} .
\]
\end{exercise}

\begin{exercise}[$\star\star$]\label{exo:g11:deriv:5}
Telaahlah variasi $f(x) = x^3 - 6x^2 + 9x - 2$ pada $\R$ (\omterm{def:g10:functions:table}{tabel
variasi}, nilai ekstrem lokal).
\end{exercise}

\begin{exercise}[$\star\star$]\label{exo:g11:deriv:6}
Telaahlah variasi $f(x) = \frac{x^2 + 3}{x + 1}$ pada
$\intoo{-1}{+\infty}$.
\end{exercise}

\begin{exercise}[$\star\star$]\label{exo:g11:deriv:7}
Misalkan $f(x) = x^2 - 4x + 5$.
\begin{enumerate}
  \item Carilah titik pada kurvanya tempat \omterm{def:g11:deriv:tangent}{garis singgungnya} mendatar.
  \item Carilah titik tempat \omterm{def:g11:deriv:tangent}{garis singgungnya} sejajar dengan garis
        $y = 2x + 1$.
\end{enumerate}
\end{exercise}

\begin{exercise}[$\star\star$]\label{exo:g11:deriv:8}
Sebuah kandang berbentuk persegi panjang dibangun bersandar pada tembok dengan
pagar sepanjang $60$ m (tiga sisi). Nyatakan luasnya sebagai \omterm{def:g11:func:function}{fungsi} lebar $x$,
lalu pakailah \omterm{def:g11:deriv:derivative}{turunan} untuk mencari ukuran yang berluas \omterm{def:g10:functions:extrema}{maksimum}. (Bandingkan
dengan metode titik puncak pada \cref{ch:g11:quad}.)
\end{exercise}

\begin{exercise}[$\star\star$]\label{exo:g11:deriv:9}
Tunjukkan bahwa untuk semua $x > 0$ berlaku $\sqrt{x} \leq \frac{x + 1}{2}$,
dengan menelaah \omterm{def:g11:func:function}{fungsi} $f(x) = \frac{x+1}{2} - \sqrt x$ pada $\intoo{0}{+\infty}$.
\end{exercise}

\begin{exercise}[$\star\star$]\label{exo:g11:deriv:10}
Sebuah kaleng berbentuk tabung harus memuat $500$ cm$^3$. Luas permukaannya
(dua cakram ditambah selimutnya) adalah $S = 2\pi r^2 + 2\pi r h$ dengan
$\pi r^2 h = 500$. Nyatakan $S$ sebagai \omterm{def:g11:func:function}{fungsi} $r$ saja, lalu tunjukkan bahwa
$S'(r) = 4\pi r - \frac{1000}{r^2}$. Simpulkan jari-jari yang meminimumkan
luas permukaannya, sebagai bentuk eksak.
\end{exercise}

\begin{exercise}[$\star\star\star$]\label{exo:g11:deriv:11}
Berapa banyak \omterm{def:g11:deriv:tangent}{garis singgung} \omterm{def:g11:quad:parabola}{parabola} $y = x^2$ yang melalui titik luar
$P(1, -3)$? Tentukan \omterm{def:g10:algebra:equation}{persamaannya}. (Misalkan $a$ \omterm{def:g10:coordgeom:system}{absis} titik
sentuhnya lalu nyatakan bahwa \omterm{def:g11:deriv:tangent}{garis singgung} di $a$ melalui $P$.)
\end{exercise}

\begin{exercise}[$\star\star\star$]\label{exo:g11:deriv:12}
Misalkan $f(x) = x^3 - 3x + 1$ (lihat \cref{ex:g11:deriv:study}). Dengan
memakai \omterm{def:g10:functions:table}{tabel variasinya}, tentukan banyaknya penyelesaian \omterm{def:g10:algebra:equation}{persamaan}
$f(x) = k$ menurut nilai \omterm{def:g10:numbers:sets}{bilangan real} $k$.
\end{exercise}

\section{Soal: Tiga pekerjaan garis singgung}

\begin{problem}[{Soal akhir pekan --- cermin parabola yang terbukti,
mesin pencari akar milik Newton, dan balok terkuat dari sebatang gelondong}]
\label{pb:g11:deriv:1}
Sebuah \omterm{def:g11:deriv:derivative}{turunan} menggambar \omterm{def:g11:deriv:tangent}{garis singgung}; kedengarannya bakat yang sederhana.
Soal ini memperlihatkan \omterm{def:g11:deriv:tangent}{garis singgung} memegang tiga pekerjaan sekaligus:
membuktikan teorema lampu sorot yang tertunda pada
\cref{pb:g11:quad:1} (setiap sinar sejajar sumbu \omterm{def:g11:quad:parabola}{parabola}
memantul melalui fokusnya), menggerakkan algoritme penyelesai
\omterm{def:g10:algebra:equation}{persamaan} tercepat yang lazim dipakai --- di dalam setiap
kalkulator dan jaringan saraf tiruan --- serta menemukan balok
persegi panjang terkuat yang dapat dihasilkan sebatang gelondong bundar.

\medskip
\noindent\textbf{Bagian I --- Kefasihan dengan \omterm{def:g11:deriv:tangent}{garis singgung}.}
\begin{enumerate}
  \item Turunkan $f(x) = 3x^2 - 5x + 1$ dan
        $g(x) = x^3 - 12x$
        (\cref{thm:g11:deriv:rules}); lalu temukan kembali $\left(x^2
        \right)'$ di $x = 1$ dari definisinya (laju
        perubahan, lalu $h \to 0$).
  \item Carilah \omterm{def:g11:deriv:tangent}{garis singgung} $y = x^3 - 12x$ di $a = 2$. Apa yang
        istimewa padanya, dan apa yang diisyaratkannya
        (\cref{prop:g11:deriv:extrema})?
  \item Tetapkan rumus umum \omterm{def:g11:deriv:tangent}{garis singgung} \omterm{def:g11:quad:parabola}{parabola} $y = x^2$
        di titik $(a, a^2)$:
        \[
        y = 2ax - a^2 .
        \]
  \item Di mana \omterm{def:g11:deriv:tangent}{garis singgung} itu memotong sumbu $y$? Simpulkan
        resep sang juru gambar: untuk menggambar \omterm{def:g11:deriv:tangent}{garis singgung} di titik $P$
        yang tingginya $h$ pada \omterm{def:g11:quad:parabola}{parabola}, hubungkan $P$ ke titik pada
        sumbunya sedalam\,\dots\ di bawah titik asal. Nyatakanlah.
  \item Susunlah \omterm{def:g10:functions:table}{tabel variasi} lengkap
        $g(x) = x^3 - 12x$ (tanda $g'$, nilai ekstrem lokal beserta
        nilainya, \cref{met:g11:deriv:study}).
\end{enumerate}

\medskip
\noindent\textbf{Bagian II --- Teorema cermin.}
Misalkan $P(a, a^2)$ sebuah titik pada \omterm{def:g11:quad:parabola}{parabola} $y = x^2$ (ambil
$a \neq 0$), $F\left(0, \frac14\right)$ fokusnya, dan $T$
titik tempat \omterm{def:g11:deriv:tangent}{garis singgung} di $P$ memotong sumbu $y$ (pertanyaan
4). Ingat kembali dari \cref{pb:g11:quad:1} bahwa
$PF = a^2 + \frac14$.
\begin{enumerate}[resume]
  \item Hitunglah jarak $FT$.
  \item Simpulkan bahwa segitiga $FPT$ sama kaki di $F$.
        Karena itu, dua sudut yang mana yang sama besar?
  \item Sinar datang yang melalui $P$, sejajar dengan sumbunya,
        sejajar pula dengan garis $(TF)$ (keduanya tegak). Dengan
        memakai sudut berseberangan, tunjukkan bahwa \omterm{def:g11:deriv:tangent}{garis singgung} di $P$
        membentuk \emph{sudut yang sama besar} dengan sinar datang yang tegak
        itu dan dengan ruas garis $[PF]$.
  \item Fisikanya: cahaya memantul pada sebuah kurva seperti pada \omterm{def:g11:deriv:tangent}{garis
        singgungnya}, dan pergi dengan sudut yang sama dengan sudut datangnya
        (hukum pemantulan). Simpulkan \emph{teorema cermin}: setiap
        sinar yang sejajar sumbunya memantul melalui fokusnya ---
        dan, dengan membalik waktunya, bola lampu di fokusnya memancarkan
        cahaya sejajar. Antena \omterm{def:g11:quad:parabola}{parabola} dan lampu sorot pada
        \cref{pb:g11:quad:1} kini menjadi teorema.
  \item Periksalah segitiga sama kaki itu secara numeris di $a = 1$:
        hitunglah $T$, $FP$ dan $FT$.
\end{enumerate}

\medskip
\noindent\textbf{Bagian III --- Mesin Newton.}
Untuk menyelesaikan $f(x) = 0$, Newton mengusulkan: berangkatlah dari terkaan
$x_0$, ikuti \omterm{def:g11:deriv:tangent}{garis singgung} di $x_0$ turun sampai sumbunya, lalu jadikan
titik potongnya dengan sumbu $x$ sebagai terkaan berikutnya.
\begin{enumerate}[resume]
  \item Turunkan rumus metodenya:
        \[
        x_{n + 1} = x_n - \frac{f(x_n)}{f'(x_n)} .
        \]
  \item Jalankan pada $f(x) = x^2 - 2$ dari $x_0 = 1$: hitunglah
        $x_1$, $x_2$, $x_3$ sebagai pecahan eksak. Barisan berumur
        dua ribu tahun yang mana --- dan mesin yang mana
        dari \cref{pb:g10:functions:1} --- yang baru saja kamu
        bangun kembali? Buktikan identitas di balik kebetulan itu:
        \[
        x - \frac{x^2 - 2}{2x} = \frac12\left(x +
        \frac2x\right).
        \]
  \item Selesaikan $x^3 = 100$ (titik penyalipan tabungan pada
        \cref{pb:g10:reffunc:1}): terapkan dua langkah Newton pada
        $f(x) = x^3 - 100$ dari $x_0 = 5$ (empat angka desimal), lalu
        bandingkan dengan $\sqrt[3]{100}$ dari kalkulator.
  \item Sabotase: cobalah menjalankan metode itu pada
        $f(x) = x^3 - 12x$ di $x_0 = 2$. Apa yang berjalan salah, dan
        mengapa (pertanyaan 2)? Nyatakan peringatan praktisnya.
  \item Pada pertanyaan 12 banyaknya angka desimal yang benar berjalan
        kira-kira $1 \to 3 \to 6$. Bandingkan dengan metode bagi dua
        (membelah dua sebuah \omterm{def:g10:numbers:interval}{interval} pada setiap langkah) lalu katakan dalam
        satu kalimat mengapa kalkulator, penerima sistem penentu posisi dan
        pelatihan jaringan saraf tiruan semuanya berjalan di atas gagasan
        Newton.
\end{enumerate}

\medskip
\noindent\textbf{Bagian IV --- Balok terkuat.}
Sebuah balok persegi panjang berlebar $w$ dan bertinggi $d$ digergaji dari
sebatang gelondong bundar berdiameter $30$ cm, sehingga $w^2 + d^2 = 900$.
Kekakuan sebuah balok sebanding dengan $w d^2$ (tingginya terhitung dua kali:
gelagar dipasang berdiri!).
\begin{enumerate}[resume]
  \item Nyatakan kekakuannya sebagai \omterm{def:g11:func:function}{fungsi} $w$ saja:
        $S(w) = 900w - w^3$, pada \omterm{def:g10:numbers:interval}{interval} yang mana?
  \item Turunkan, susun \omterm{def:g10:functions:table}{tabel variasinya}, lalu carilah
        $w$ dan $d$ yang optimum (nilai eksaknya, lalu sampai satuan
        milimeter).
  \item Tunjukkan bahwa perbandingan optimumnya memenuhi
        $d = w\sqrt2$ --- rasio para tukang kayu zaman dahulu. (Cerita
        tambahan: membagi diameter gelondongnya menjadi tiga bagian sama
        panjang lalu menegakkan garis tegak lurus sampai lingkarannya
        melukis persis persegi panjang ini.)
  \item Pilihan yang naif adalah balok \emph{persegi}
        ($w = d$). Hitunglah kekakuannya dan kekakuan balok
        optimumnya, lalu berikan keuntungan tukang kayunya dalam persen.
  \item Penutup --- ketiga pekerjaan \omterm{def:g11:deriv:tangent}{garis singgung} itu, masing-masing
        dalam satu kalimat:
        geometri (teorema cermin), numerik (mesin
        Newton), pengoptimuman (baloknya), semuanya digerakkan oleh
        objek yang sama. Lalu penunjuk ke depan: tahun berikutnya menambahkan
        kecembungan --- sisi kurva tempat \omterm{def:g11:deriv:tangent}{garis singgungnya} berada --- untuk
        mempertajam ketiganya.

\end{enumerate}
\end{problem}
